\newpage
\section{Resultater}
Dette afsnit viser, hvordan metoden ser ud, når den anvendes på Onsdagsbanen 2026. Først vises de kontrafaktiske tider løb for løb. Derefter følger et kort eksempel, hvor estimerede og observerede tider sammenholdes for to konkrete både.

\subsection*{Kontrafaktiske sejltider ved fravær}
Tabellerne nedenfor viser 2026-løbene hver for sig. For hvert løb vises sejltid, beregnet tid, kontrafaktisk sejltid og kontrafaktisk beregnet tid.
Både, som ikke deltog i et løb, har tomme felter for sejltid, beregnet tid og afvigelser. Den kontrafaktiske sejltid er markeret med gult for disse både.
For både, som faktisk deltog i løbet, skal kolonnerne Est. Tid og Est. Hdcp. Tid læses som one-step leave-one-out-forudsigelser. Det vil sige, at bådens egen observerede tid ikke indgår i dagens fit, hverken ved estimeringen af dagens gns. sejltid eller dagens kaosniveau. For både, som ikke deltog, er de estimerede tider derimod rene kontrafaktiske forudsigelser baseret på bådens prior og dagens øvrige felt.

\input{\MissingRacePredictionTablesTex}

\subsection{Sammenligninger mellem et par både}
I figuren herunder ses hvordan Rikkes og Caspers niveau udvikler sig over tid.
Det er interessant hvordan dage med meget kaos resulterer i en a priori prediktion som flytter sig i forhold til det kaos, alt efter om den pågældende båd har negativt eller postitivt $k_{i,t}$. Det generelle niveau ses tydeligere i ræs med mindre kaos. Man kan se at usikkerheden (en standardafvigelse) på det underliggende niveau generelt ikke er så høj, men at variansen på dagen spiller en stor rolle. Rikkes forudsagte respons på kaos går fra at være gunstigt, til ugunstigt og tilbage igen. Caspers niveau er estimeret til at være generelt bedre, og hans respons på kaos er også bedre i gennemsnit.

\DualBoatPercentPlotFromCsv{%
title={Casper Lyhne og Rikke M. Schnuchel},
first-name={C. Lyhne},
first-csv={\BoatPlotDataDir/boat_plot_casper-lyhne_percent.csv},
first-color={green!70!black},
second-name={R. M. Schnuchel},
second-csv={\BoatPlotDataDir/boat_plot_rikke-m-schnuchel_percent.csv},
second-color={magenta},
ymin={-16.090455255458},
ymax={12.612653426958}%
}

På dette plot ser vi Axel Thomsen og Casper Lyhne sammen. Axel har et ekstremt godt respons på kaos, og har også generel bedre fart end Casper. Han sejler ofte og konsistent, så usikkerheden på hans underliggende niveau og kaos respons er lav i forhold til Caspers.

\DualBoatPercentPlotFromCsv{%
title={Axel Thomsen og Casper Lyhne},
first-name={C. Lyhne},
first-csv={\BoatPlotDataDir/boat_plot_casper-lyhne_percent.csv},
first-color={green!70!black},
second-name={A. Thomsen},
second-csv={\BoatPlotDataDir/boat_plot_axel-thomsen_percent.csv},
second-color={blue!70!black},
ymin={-20.250542179598},
ymax={12.064646902344}%
}

På dette plot sammenlignes Christian Grejs og Axel Thomsen. De har begge ekstremt god kaos respons, men Axels er bedst. Det er dog ikke nok til at kompensere for ren boat speed hos Christian.

\DualBoatPercentPlotFromCsv{%
title={Christian Grejs og Axel Thomsen},
first-name={C. Grejs},
first-csv={\BoatPlotDataDir/boat_plot_christian-grejs_percent.csv},
first-color={red},
second-name={A. Thomsen},
second-csv={\BoatPlotDataDir/boat_plot_axel-thomsen_percent.csv},
second-color={blue!70!black},
ymin={-25.5},
ymax={13.0}%
}

\subsection*{Trajektorier i $(s,k)$ på Stor Bane}
Her ser vi en samling af alles generelle niveau og kaos respons.
Det viser sig at der er en meget tydelig korrelation mellem at være generelt hurtig og respondere godt på kaos. Men der er mange både der stort set ingen kaos respons har - hverken negativ eller positiv. Lysere farver indikerer ældre data.

\XGammaTrajectoryPlotFromCsv{}

\input{\StorBaneSplitRecommendationTex}
